\documentclass[10pt,a4paper]{article}

\usepackage[margin=1in]{geometry}
\usepackage{setspace}
\setstretch{1.15}
\usepackage[UTF8,fontset=fandol]{ctex}
\usepackage{graphicx}
\usepackage{amssymb,amsthm,amsmath}
\usepackage{xcolor}
\usepackage{hyperref}
\usepackage{booktabs}
\usepackage{array}
\usepackage{longtable}
\usepackage{tabularx}
\usepackage{enumitem}
\usepackage[section]{placeins}

\hypersetup{
    colorlinks=true,
    linkcolor=blue,
    filecolor=blue,
    urlcolor=blue,
    citecolor=blue
}

\renewcommand{\topfraction}{0.9}
\renewcommand{\bottomfraction}{0.8}
\renewcommand{\textfraction}{0.07}
\renewcommand{\floatpagefraction}{0.8}
\setcounter{topnumber}{3}
\setcounter{bottomnumber}{2}
\setcounter{totalnumber}{5}

\newcommand{\otf}{\textsf{OTF}}
\newcommand{\tp}{T_{pp}}
\newcommand{\ts}{T_{ss}}
\newcommand{\lam}{\lambda}
\newcommand{\ang}{\theta}
\newcommand{\loss}{\mathcal{L}}
\newcommand{\Rtar}{\mathbf{R}^{\star}}
\newcommand{\Rpred}{\widehat{\mathbf{R}}}

\newcommand{\placeholderfigure}[2]{%
\begin{figure}[!htbp]
\centering
\fbox{%
\begin{minipage}[c][2.1in][c]{0.9\linewidth}
\centering
\textbf{Figure Placeholder}\\[0.4em]
\small #1
\end{minipage}}
\caption{#2}
\end{figure}
}

\begin{document}

\title{补充材料\\[0.4em]\large 物理一致的角分辨傅里叶算子超表面生成式逆向设计}
\author{}
\date{\today}
\maketitle

本补充材料用于展开正文中被简写的方法细节。行文方式将尽量保持论文补充材料风格，而不是工程说明风格，重点围绕数据集构建、前向代理模型，以及条件扩散逆向设计的理论与结构。

\section{数据集制备与任务打分}

\subsection{自由形态结构生成与物理标注}

本文考虑的设计空间是尺寸为 $64\times64$ 的自由形态二值超表面图案，其中 $1$ 表示材料区域，$0$ 表示空气区域。每个结构由粗网格随机场出发，经 $C_4+\sigma_x$ 对称约束、平滑、二值化和形态学规整后得到。这样的构造方式一方面显著扩展了可探索的几何形态，另一方面又避免了过细桥连、小孤岛和微小孔洞等不利于制造的局部细节。

数据集并不对应单一填充率，而是在
\begin{equation}
\mathrm{fill}\in\{0.4,0.5,0.6,0.7,0.8\}
\end{equation}
五个名义填充率档位上进行分布，同时进一步扰动粗网格尺度、平滑强度和最小特征阈值。这样得到的结构集合既包含拓扑层面的自由度，也包含材料占空比层面的系统变化。

每个结构都通过 RCWA 进行标注。本文采用固定物理平台：硅图形层位于二氧化硅基底上，入射介质为空气；周期为 $500$~nm，结构层厚度为 $500$~nm。响应在波长、入射角和偏振保持传输通道上离散采样，因此训练目标不是单一器件标签，而是一个离散响应张量。

\begin{table}[!htbp]
\centering
\caption{数据集与物理标注配置。}
\begin{tabularx}{\linewidth}{>{\raggedright\arraybackslash}p{0.36\linewidth}X}
\toprule
项目 & 配置 \\
\midrule
结构表示 & 自由形态二值图案，$64\times64$ 像素 \\
对称先验 & $C_4+\sigma_x$ \\
二值约定 & $1=\mathrm{material}$，$0=\mathrm{air}$ \\
名义填充率档位 & $0.4, 0.5, 0.6, 0.7, 0.8$ \\
周期 & $500$ nm \\
结构层厚度 & $500$ nm \\
结构材料 & Si \\
入射 / 输出介质 & Air / SiO$_2$ \\
波长采样 & $800$--$1300$ nm，步长 $50$ nm（共 $11$ 点） \\
角度采样 & $-40^\circ$ 到 $40^\circ$，步长 $5^\circ$（共 $17$ 点） \\
保留通道 & $\tp$ 与 $\ts$ 幅值 \\
响应张量尺寸 & $2\times11\times17$ \\
\bottomrule
\end{tabularx}
\end{table}

\subsection{响应张量表示}

对每个结构，物理标签可写为
\begin{equation}
\mathbf{R}(\lam,\ang)=
\begin{bmatrix}
\tp(\lam,\ang)\\
\ts(\lam,\ang)
\end{bmatrix}
\in \mathbb{R}^{2\times 11\times 17}.
\end{equation}
因此，学习框架始终围绕两个对象展开：结构图像 $x\in\{0,1\}^{64\times64}$ 和响应张量 $\mathbf{R}$。在训练阶段，对响应张量采用逐格点标准化：
\begin{equation}
\widetilde{\mathbf{R}}_{qjk}=\frac{\mathbf{R}_{qjk}-\mu_{qjk}}{\sigma_{qjk}+\varepsilon},
\end{equation}
其中 $q$ 表示偏振通道，$j$ 表示波长索引，$k$ 表示角度索引。这样的处理不会把零点、泄漏带和弱响应肩部等物理上重要但数值较小的位置平均掉。

\subsection{任务打分函数族}

本文并不使用单一通用分数来覆盖全部算子目标，而是按物理功能定义一组 task-aware 指标。对于一维行算子目标，理想角谱包络写为
\begin{equation}
g_m(\theta)=\left|\frac{\sin\theta}{\sin\theta_{\max}}\right|^m,
\end{equation}
其中 $m=2,3,4$ 分别对应二阶、三阶和四阶算子。每个波长切片上的总分由中心抑制、形状拟合、边缘支撑和带宽保持共同构成：
\begin{equation}
S_m=(1-w_b)\left(w_cS_c+w_sS_s+w_eS_e\right)+w_bS_b.
\end{equation}
本文当前采用
\begin{equation}
w_c=0.6,\qquad w_s=0.3,\qquad w_e=0.1,\qquad w_b=0.2.
\end{equation}

对偏振无关任务，两个通道都必须保持目标行响应，并在大角度处彼此匹配：
\begin{align}
S_{\mathrm{PI}}&=(1-\beta)\min(S_p,S_s)+\beta S_{\mathrm{match}},\\
S_{\mathrm{match}}&=1-\frac{|T_{pp}(+\theta_a)-T_{ss}(+\theta_a)|+|T_{pp}(-\theta_a)-T_{ss}(-\theta_a)|}{2A_{\max}},
\end{align}
其中 $\beta=0.15$，$\theta_a=40^\circ$。

对偏振复用任务，一个通道负责执行目标，另一个通道尽量静默：
\begin{align}
S_{\mathrm{sil}}&=\max\left(0,1-\frac{R_{\mathrm{passive,max}}}{\tau}\right),\\
S_{\mathrm{PM}}&=\frac{w_aS_{\mathrm{act}}+w_qS_{\mathrm{sil}}}{w_a+w_q},
\end{align}
当前取 $w_a=0.7$、$w_q=0.3$。

对低通目标，理想角谱包络写为
\begin{equation}
g_\sigma(\theta)=\exp\!\left(-\frac{\theta^2}{\sigma^2}\right),
\end{equation}
并在 $\sigma\in\{8^\circ,12^\circ,16^\circ\}$ 的候选集合中搜索最合适的等效通带宽度。当前排序主要由形状项驱动，而中心通带保持和边缘抑制则作为诊断量保留。

对时空微分任务，目标定义在二维 wavelength--theta 窗口上：
\begin{equation}
G^{\star}(\lambda,\theta)\propto k_x^2(\theta)\bigl(\omega(\lambda)-\omega_0\bigr)^2.
\end{equation}
其总分写为
\begin{equation}
S_{\mathrm{ST}}=0.50S_{\mathrm{zero}}+0.40S_{\mathrm{corner}}+0.10S_{\mathrm{aux}},
\end{equation}
其中辅助项同时考虑加权模板误差与模式投影一致性。

\begin{table}[!htbp]
\centering
\caption{本文使用的任务打分摘要。}
\begin{tabularx}{\linewidth}{>{\raggedright\arraybackslash}p{0.18\linewidth}>{\raggedright\arraybackslash}p{0.20\linewidth}>{\raggedright\arraybackslash}p{0.28\linewidth}X}
\toprule
任务 & 目标形式 & 主要分数组成 & 物理含义 \\
\midrule
二阶 / 三阶 / 四阶 & 行包络 $g_m(\theta)$ & 中心、形状、边缘、带宽 & 零点、增长规律、外角响应、有限带宽保持 \\
偏振无关 & 双通道共享行响应 & 双通道行分数 + 大角度匹配 & 同时实现目标功能并保持偏振一致性 \\
偏振复用 & 主动行 + 静默被动行 & 主动通道形状 + 被动抑制 & 一个通道工作、另一个通道安静 \\
低通 & 高斯行包络 $g_\sigma(\theta)$ & 形状主导，辅以中心 / 边缘诊断 & 中心通带与外角抑制 \\
时空微分 & 2D 窗口 $k_x^2(\omega-\omega_0)^2$ & 零线、角落、模板误差、模式投影 & 联合波长--角度形貌而非单行匹配 \\
\bottomrule
\end{tabularx}
\end{table}

\placeholderfigure
{建议后续正式配图：自由形态结构生成、RCWA 标注网格、行打分与二维窗口打分的统一示意。}
{数据集制备与任务打分占位图。正式图可将结构生成、物理标注和典型评分分解统一到一张总览图中。}

\section{前向代理模型结构与设计考虑}

\subsection{前向映射目标}

前向代理模型所近似的是
\begin{equation}
\mathcal{F}_{\phi}: \{0,1\}^{64\times64}\rightarrow \mathbb{R}^{2\times11\times17},
\end{equation}
即从结构直接预测整个双偏振 wavelength--theta 响应张量。它的作用边界很明确：不直接替代严格电磁仿真给出最终结论，而是为候选筛选和逆向训练提供快速的响应近似。

\subsection{结构设计}

当前 surrogate 采用带坐标注入的卷积编码器与多尺度响应网格融合结构。输入端首先将二值结构图与两个笛卡尔坐标通道拼接，使模型能够把平面布局与角分辨散射方向性联系起来。随后，多层残差编码器在不同尺度上抽取局部细节与全局拓扑特征。最深层 latent 经过全局池化和 MLP 后广播回整个响应网格，使每个输出格点都能够感知填充率、连通性与整体拓扑信息。最后，各尺度特征在 $11\times17$ 的响应平面上融合，并与显式的响应网格坐标共同输入输出头，得到双通道预测。

因此，整体结构可概括为
\begin{equation}
\Rpred = H\!\left(P_0,P_1,P_2,P_3,G_{\mathrm{global}},C_{\lambda\theta}\right),
\end{equation}
其中 $P_\ell$ 表示不同尺度特征在响应网格上的投影，$G_{\mathrm{global}}$ 表示全局 latent 广播分支，$C_{\lambda\theta}$ 表示响应平面的显式坐标场。

\subsection{物理考虑}

输入坐标注入的必要性在于，角分辨散射不是普通的平移不变纹理问题。多尺度编码是必要的，因为器件性能同时依赖边界级局部结构和整体拓扑。全局广播分支则使每个波长--角度采样点都能获得长程结构信息。两个偏振通道采用联合输出，也更符合物理对象本身的性质，因为 $\tp$ 和 $\ts$ 都由同一几何产生，其相关性本身就是偏振无关和偏振复用任务的重要组成部分。

\subsection{训练设置}

前向代理训练使用标准化后的响应张量作为监督信号，主损失为
\begin{equation}
\loss_{\mathrm{fwd}}=\left\|\mathcal{F}_{\phi}(x)-\widetilde{\mathbf{R}}\right\|_1.
\end{equation}
当前训练配置采用 batch size $64$、AdamW 优化器、学习率 $10^{-4}$、权重衰减 $10^{-4}$、$90\%/10\%$ 的训练/验证划分、梯度裁剪、学习率自适应衰减与 early stopping。训练时同时使用保持物理定义不变的对称翻转增强。

\begin{table}[!htbp]
\centering
\caption{前向代理模型的详细架构。}
\begin{tabularx}{\linewidth}{>{\raggedright\arraybackslash}p{0.18\linewidth}>{\raggedright\arraybackslash}p{0.18\linewidth}>{\raggedright\arraybackslash}p{0.17\linewidth}X}
\toprule
阶段 & 分辨率 / 通道 & 主要操作 & 架构作用 \\
\midrule
输入组装 & $64\times64$，$1+2$ 通道 & 二值结构与 $(x,y)$ 坐标拼接 & 打破纯平移不变性并显式提供几何位置信息 \\
Stem block & $64\times64$，$32$ 通道 & 卷积 + 残差细化 & 提取低层边界与填充特征 \\
编码阶段 1 & $32\times32$，$64$ 通道 & 残差下采样 & 捕捉短程形貌并扩大感受野 \\
编码阶段 2 & $16\times16$，$128$ 通道 & 残差下采样 & 融合中尺度拓扑与局部结构母题 \\
编码阶段 3 & $8\times8$，$256$ 通道 & 残差下采样 & 聚合更强的全局结构信息 \\
latent 精炼 & $8\times8$，$256$ 通道 & 堆叠残差块 & 稳定最深层高阶表征 \\
全局上下文分支 & 池化 latent $\rightarrow128$ 通道 & 自适应池化 + MLP + 广播 & 向全部输出位置注入整体拓扑与有效填充信息 \\
多尺度投影 & 各尺度 $\rightarrow 11\times17$ & 双线性插值到响应网格 & 让全部特征在物理采样平面对齐 \\
网格坐标注入 & $11\times17$，$2$ 通道 & 显式响应平面坐标 & 区分中心 / 外角与短波 / 长波位置 \\
融合模块 & $11\times17$，拼接后的多尺度张量 & 卷积融合 & 统一局部、全局和坐标信息 \\
输出头 & $11\times17$，$2$ 通道 & 最终卷积预测 & 联合输出 $\tp$ 与 $\ts$ 响应图 \\
\bottomrule
\end{tabularx}
\end{table}

\begin{table}[!htbp]
\centering
\caption{前向代理模型训练协议摘要。}
\begin{tabularx}{\linewidth}{>{\raggedright\arraybackslash}p{0.28\linewidth}X}
\toprule
项目 & 设置 \\
\midrule
损失函数 & 标准化响应空间中的逐格点 $L_1$ 回归 \\
优化器 & AdamW \\
学习率 / 权重衰减 & $10^{-4}$ / $10^{-4}$ \\
Batch size & $64$ \\
训练 / 验证划分 & $90\%$ / $10\%$ \\
稳定化策略 & 梯度裁剪、学习率自适应衰减、early stopping \\
数据增强 & 仅采用保持物理定义不变的对称翻转 \\
验证读出 & 同时关注标准化损失与反标准化后的物理空间 MAE \\
\bottomrule
\end{tabularx}
\end{table}

\placeholderfigure
{建议后续正式配图：输入结构与坐标通道、残差编码层级、全局广播、多尺度响应网格融合、双通道输出头。}
{前向代理模型占位图。正式图可将坐标感知输入、多尺度编码与响应网格融合画成一张清晰的架构图。}

\section{条件扩散逆向设计：理论、结构与条件嵌入}

\subsection{扩散正向过程}

由于同一目标响应往往对应多个可能结构，逆向问题被表述为条件扩散生成过程。设 $x_0\in[-1,1]^{1\times64\times64}$ 表示由二值结构线性映射得到的干净样本，则正向加噪马尔可夫链可写为
\begin{equation}
q(x_t|x_{t-1})=\mathcal{N}\!\left(x_t;\sqrt{1-\beta_t}\,x_{t-1},\beta_t I\right),
\end{equation}
其中 $\beta_t$ 由 cosine schedule 给出。经多步递推后，$x_t$ 关于 $x_0$ 的边缘分布可写成闭式：
\begin{equation}
q(x_t|x_0)=\mathcal{N}\!\left(x_t;\sqrt{\bar{\alpha}_t}\,x_0,(1-\bar{\alpha}_t)I\right),
\end{equation}
其中
\begin{equation}
\alpha_t=1-\beta_t,\qquad \bar{\alpha}_t=\prod_{s=1}^{t}\alpha_s.
\end{equation}
于是可以得到熟悉的重参数化形式：
\begin{equation}
x_t=\sqrt{\bar{\alpha}_t}\,x_0+\sqrt{1-\bar{\alpha}_t}\,\epsilon,\qquad \epsilon\sim\mathcal{N}(0,I).
\end{equation}

\subsection{反向过程与 velocity 参数化推导}

目标是在给定条件 $\mathbf{c}$ 的情况下学习反向过程
\begin{equation}
p_{\theta}(x_{t-1}|x_t,\mathbf{c}).
\end{equation}
在扩散模型中，常见的预测对象包括噪声 $\epsilon$、干净样本 $x_0$，或二者的线性组合。本文采用的是 velocity 参数化，其目标定义为
\begin{equation}
v^\star=\sqrt{\bar{\alpha}_t}\,\epsilon-\sqrt{1-\bar{\alpha}_t}\,x_0.
\end{equation}
之所以采用这种形式，是因为它和 $(x_t,x_0,\epsilon)$ 之间具有简洁的线性可逆关系。由
\begin{align}
x_t &= \sqrt{\bar{\alpha}_t}\,x_0+\sqrt{1-\bar{\alpha}_t}\,\epsilon,\\
v^\star &= \sqrt{\bar{\alpha}_t}\,\epsilon-\sqrt{1-\bar{\alpha}_t}\,x_0,
\end{align}
可直接解出
\begin{align}
\hat{x}_0 &= \sqrt{\bar{\alpha}_t}\,x_t-\sqrt{1-\bar{\alpha}_t}\,\hat{v},\\
\hat{\epsilon} &= \sqrt{1-\bar{\alpha}_t}\,x_t+\sqrt{\bar{\alpha}_t}\,\hat{v},
\end{align}
其中 $\hat{v}=v_{\theta}(x_t,t,\mathbf{c})$ 是网络预测。也就是说，一旦预测了 velocity，就能同时恢复去噪结构估计和噪声估计，这也是当前参数化在数值上较稳定的原因之一。

\subsection{训练目标与物理约束}

扩散基本损失为 velocity 空间中的均方误差：
\begin{equation}
\loss_{\mathrm{diff}}=\mathbb{E}_{x_0,\epsilon,t}\left[\left\|v_{\theta}(x_t,t,\mathbf{c})-v^\star\right\|_2^2\right].
\end{equation}
但仅靠扩散重建损失并不足以保证生成结果具有正确的电磁响应，因此总目标写为
\begin{equation}
\loss_{\mathrm{total}}=\loss_{\mathrm{diff}}+\lambda_{\mathrm{phys}}\loss_{\mathrm{phys}}+\lambda_{\mathrm{bin}}\loss_{\mathrm{bin}}.
\end{equation}
其中响应一致性项为
\begin{equation}
\loss_{\mathrm{phys}}=\left\|\mathcal{F}_{\phi}\!\left(\frac{\hat{x}_0+1}{2}\right)-\mathbf{c}\right\|_1,
\end{equation}
即利用前向代理模型比较去噪样本的预测响应与目标条件；二值正则项为
\begin{equation}
\loss_{\mathrm{bin}}=\mathrm{mean}\!\left(x(1-x)\right),
\end{equation}
用于压制灰区不确定性。本文当前使用
\begin{equation}
\lambda_{\mathrm{phys}}=0.5,\qquad \lambda_{\mathrm{bin}}=0.05.
\end{equation}

\subsection{条件表示与条件嵌入方式}

模型条件是双通道响应张量
\begin{equation}
\mathbf{c}\in\mathbb{R}^{2\times11\times17}.
\end{equation}
条件编码器并不只输出单一向量，而是同时产生三类表示：
\begin{enumerate}[leftmargin=1.5em]
\item 全局条件向量；
\item 多尺度条件特征图；
\item spectral token 序列。
\end{enumerate}
它们通过三条不同路径进入 U-Net：
\begin{enumerate}[leftmargin=1.5em]
\item \textbf{条件图注入：} 多尺度条件图经过尺寸匹配后，以加性方式注入残差块；
\item \textbf{全局条件调制：} 全局条件向量与时间嵌入共同生成 scale--shift 参数，用来调制残差块内部的归一化特征；
\item \textbf{token 级 cross-attention：} 深层 latent 特征通过 cross-attention 直接访问条件 token，从而耦合空间生成与全局光谱条件。
\end{enumerate}

\subsection{classifier-free guidance 与采样}

训练时，模型以一定概率丢弃真实条件并用可学习的空条件替代，从而实现 classifier-free guidance 的训练准备。推理时，conditional 与 unconditional 的预测按
\begin{equation}
v_{\mathrm{cfg}}=v_{\mathrm{uncond}}+s_{\mathrm{cfg}}\left(v_{\mathrm{cond}}-v_{\mathrm{uncond}}\right)
\end{equation}
进行组合，其中默认引导强度为 $s_{\mathrm{cfg}}=3.0$。采样从高斯噪声出发，沿时间步逐步反推回 $x_0$，再从 $[-1,1]$ 映射回 $[0,1]$ 并以阈值 $0.5$ 二值化，最终得到候选结构。

\begin{table}[!htbp]
\centering
\caption{条件扩散逆向模型的详细架构。}
\begin{tabularx}{\linewidth}{>{\raggedright\arraybackslash}p{0.20\linewidth}>{\raggedright\arraybackslash}p{0.18\linewidth}>{\raggedright\arraybackslash}p{0.18\linewidth}X}
\toprule
模块 & 分辨率 / 形式 & 主要操作 & 架构作用 \\
\midrule
加噪输入支路 & $64\times64$，$1$ 通道 & 扩散状态输入卷积 & 初始化空间去噪主干 \\
时间嵌入路径 & 向量嵌入 & 正弦时间嵌入 + MLP & 为全部残差块编码扩散步信息 \\
条件编码器 stem & $11\times17$，多通道输入 & 带坐标的卷积编码 & 从目标响应张量抽取结构化条件特征 \\
条件下采样路径 & $11\times17 \rightarrow 6\times9 \rightarrow 3\times5$ & 残差下采样 & 构建与 U-Net 深度对齐的多尺度条件图 \\
全局条件向量 & 紧凑向量表征 & 池化后的条件摘要 & 为 scale--shift 调制提供全局响应上下文 \\
光谱 token 投影 & token 序列 & 展平 + 线性投影 & 将深层条件特征转换为可注意力访问的 token \\
编码阶段 1 & 浅层空间尺度 & 残差块 + 条件图注入 & 在高分辨率阶段引入局部条件约束 \\
编码阶段 2 & 中间空间尺度 & 残差块 + 条件图注入 & 将中尺度几何演化与中尺度条件图绑定 \\
深层编码阶段 & 深层空间尺度 & 残差块 + self-attention + cross-attention & 让 latent 几何特征访问全局条件 token \\
瓶颈模块 & 最深 latent 层级 & 残差块 + self-attention + cross-attention & 在最压缩表征中进行全局条件化去噪 \\
解码阶段 1 & 深层上采样尺度 & skip 融合 + 残差块 + cross-attention & 在重建初期保持 token 级全局条件约束 \\
解码阶段 2 & 中间上采样尺度 & skip 融合 + 残差块 & 恢复中尺度几何结构并承接条件图信息 \\
解码阶段 3 & 浅层上采样尺度 & skip 融合 + 残差块 & 重建高分辨率边界与二值轮廓 \\
输出投影 & $64\times64$，$1$ 通道 & 归一化 + 最终卷积 & 在结构平面上预测 velocity 场 \\
Classifier-free guidance & conditional / unconditional 对 & 采样时的引导组合 & 强化反向扩散中的目标条件约束 \\
Physics consistency 分支 & surrogate 评估的去噪样本 & 响应一致性正则 & 将生成过程重新绑定到目标电磁响应 \\
二值正则 & 映射到 $[0,1]$ 的去噪样本 & 灰值惩罚 & 鼓励输出接近清晰二值结构 \\
\bottomrule
\end{tabularx}
\end{table}

\placeholderfigure
{建议后续正式配图：扩散正向过程、velocity 参数化、条件 U-Net、三类条件嵌入路径与 classifier-free guidance。}
{条件扩散逆向设计占位图。正式图建议将扩散公式推导、U-Net 主干以及条件图 / 全局向量 / token 三条条件嵌入路径统一展示。}

\section{对比模型与评估维度}

\subsection{对比模型}

本文在统一的 response-centered 设定下比较几类具有代表性的逆向设计范式。Topology optimization 直接围绕目标任务分数迭代更新结构，其优点是物理指向明确，但当目标种类增多或需要探索多个候选解时，计算代价会迅速上升。

CVAE 将目标响应作为条件，通过连续潜变量生成结构，优点是训练与采样都较稳定，但容易把本来多模态的可行解集压缩成更平滑、更平均的结果。cGAN 同样执行从目标响应到结构的条件生成，只是以对抗训练为核心。它可以生成更锐利的结构形态，但训练稳定性和对多种可行模式的覆盖通常更难兼顾。

前向引导扩散模型与这些基线的差别主要有两点。第一，反向去噪过程天然更适合描述一对多的条件分布，因此与超表面逆向设计问题本身更匹配。第二，前向代理在训练中以响应一致性约束的方式参与，使生成过程不仅受到数据分布约束，还持续受到目标电磁响应的牵引。

\subsection{评估维度}

所有对比模型都在相同目标条件下，并采用同一套仿真后验证逻辑进行评估。本文把比较维度分成三类：
\begin{enumerate}[leftmargin=1.5em]
\item \textbf{任务表现：} best score、mean score、top-$k$ score 和阈值成功率；
\item \textbf{响应保真：} 全波验证后的 response / spectrum MAE；
\item \textbf{生成质量：} binary rate、diversity 和 inference time。
\end{enumerate}
这些维度分别衡量模型能否稳定命中高质量解、其真实光学响应是否贴近目标，以及生成结果是否适合后续制造导向筛选。

\begin{table}[!htbp]
\centering
\caption{本文对比模型及评估维度摘要。}
\begin{tabularx}{\linewidth}{>{\raggedright\arraybackslash}p{0.23\linewidth}>{\raggedright\arraybackslash}p{0.31\linewidth}X}
\toprule
模型 / 维度 & 核心思想 & 主要揭示的问题 \\
\midrule
Topology optimization & 围绕任务目标进行迭代结构更新 & 物理直接，但更偏向高代价单目标搜索 \\
CVAE & 条件潜变量生成 & 采样高效，但潜空间压缩可能削弱多模态覆盖 \\
cGAN & 对抗式条件生成 & 输出较锐利，但模式覆盖与训练稳定性更难兼顾 \\
前向引导 diffusion & 带响应一致性约束的迭代去噪生成 & 更适合一对多逆向设计并朝高价值响应区域收缩 \\
\midrule
任务表现 & best、mean、top-$k$、success rate & 衡量模型能否反复找到高质量验证解 \\
响应保真 & 验证后的 response / spectrum MAE & 衡量真实光学响应是否逼近目标条件 \\
生成质量 & binary rate、diversity、inference time & 衡量输出是否适合制造、是否多样、以及获得成本 \\
\bottomrule
\end{tabularx}
\end{table}

\subsection{对比分析}

前向引导扩散模型的优势，并不只是生成出的结构在视觉上更像二值超表面，而是它与逆向问题本身的结构更匹配。对同一个目标响应而言，可行结构往往天然是一对多的，因此扩散模型比单次确定性回归更适合去近似整个条件可行集。

前向代理进一步强化了这一优势。如果没有响应引导，生成模型即便贴近训练数据分布，也可能偏离真正目标传递函数所对应的物理解空间。而在本文中，扩散模型在训练过程中持续通过前向代理检查去噪候选是否仍然朝目标响应靠近，因此更不容易偏离高价值解区域。对于本文这类“高分解稀疏但真实存在”的任务，这一点尤其重要。

因此，前向引导扩散模型更有能力同时兼顾三件事：较高的验证任务分数、较低的响应失配，以及足够有用的候选多样性。这也是它成为本文核心逆向设计模型的主要原因。

\placeholderfigure
{建议后续正式配图：Topology optimization、CVAE、cGAN 与前向引导 diffusion 在任务表现、响应保真与生成质量三个维度上的统一比较。}
{对比模型总览占位图。正式图可将代表性逆向设计范式与统一评估维度放在同一张图中展示。}

\end{document}
